\section*{Introduction}

Optimisation problems are of the following form.

\begin{problem}[Optimisation Problem]
    We want to minimise \(f\pqty{\vb{x}}\), subject to some constraints \(\vb{h}\pqty{\vb{x}} = \vb{b}\), for \(\vb{x} \in \mathcal{X}\).
\end{problem}

We may have some inequality constraint instead of an equality, with \(\vb{h}\pqty{\vb{x}} \leq \vb{b}\), and we impose this component-wise in higher dimensions.

In fact, inequalities constraints can be transformed into equality constraints, as
\[
    \vb{h}\pqty{\vb{x}} \leq \vb{b}
\]
is the same as
\[
    \vb{h}\pqty{\vb{x}} + \vb{s} = \vb{b}
\]
for \(\vb{s} \geq \vb{0}\), and we also minimise over \(\vb{s}\).

\begin{definition}[Terminology on Optimisation]
    \label{def:terminology}%
    The function \(\functiondomain{f}{\RR^n}{\RR}\) that we are trying to minimise is the \emph{objective function}; \(\vb{h}\pqty{\vb{x}} = \vb{b}\) where \(\functiondomain{h}{\RR^n}{\RR^m}\) is a \emph{functional constraint} with \(\vb{b} \in \RR^m\); and \(\mathcal{X}\) (often being \(\vb{x}\) being in the positive quadrant, or a cube/ball), is a \emph{regional constraint}.

    The \(\vb{x}^*\) which solves the problem is called the \emph{optimal solution}, and the value of the function at \(\vb{x}^*\), \(f\pqty{\vb{x}^*}\), is called the \emph{optimal value/cost}.

    The set
    \[
        \mathcal{X}\pqty{\vb{b}} = \set{x \in \mathcal{X}}{\vb{h}\pqty{\vb{x}} = \vb{b}}
    \]
    is called the \emph{feasible set}.

    The extra variables \(s\) are called \emph{slack variables}.
\end{definition}
